\answer{ex:16:1}{혈액은 $\tau=t_{1/2}/\ln2=14.4$~min인 RC 회로다. 최댓값 대 최솟값은 $2^{T/t_{1/2}}=64$; 기본 고조파는 $0.55$로 감쇠한다. 비가 $2$가 되는 것은 $t_{1/2}=T=60$~min일 때다.}
\answer{ex:16:2}{$\varphi^*=\arcsin\bigl((\omega-\omega_0)/K_\varphi\bigr)\approx41$~min ($\omega-\omega_0\approx0.047$~rad/day), 이완 시간 $1/(K_\varphi\cos\varphi^*)\approx3.9$일; $+6$~h와 $-6$~h 이동 뒤에는 각각 $7.7$일과 $7.9$일.}
\answer{ex:16:3}{$K_c(n)=\sec^n(\pi/n)$: $8$과 $2.37$, 오차 $11\%$와 $30\%$; $n\to\infty$에서 $K_c\to1$, 오차 $\to50\%$.}
\answer{ex:16:4}{극단 명령(가속 페달 $80$, 브레이크 $0$)으로 $f=120$에 $0.76$~s 만에 도달한다; “가속 페달만”이면 $2.33$~s로 세 배 걸린다.}
\answer{ex:16:5}{$3\arctan(\omega\tau)+\omega d=\pi$, $K_c=(1+\omega^2\tau^2)^{3/2}$: $K_c=8$, $3.54$, $2.50$, $1.76$, 주기 $3.63$, $5.46$, $6.86$, $9.27\,\tau$; $0.9K_c$에서는 진동이 감쇠하고 $1.1K_c$에서는 일정한 진동이 자리 잡는다.}
\answer{ex:16:6}{$0.60\bar c$, $0.87\bar c$, $0.90\bar c$, 극한 $0.91\bar c$; 농도가 일정하면 반응은 0이다. 이런 수신자는 분량만 읽는다.}
\answer{ex:16:7}{열린 문제. 되먹임이라면 리듬은 $K>8$에서만 있고, 주기는 단의 $\tau$에 비례하며 $K$에 거의 의존하지 않는다($K=8.5$에서 $3.64\,\tau$, $K=40$에서 $4.23\,\tau$). $K$를 1.5배 바꾸면 주기는 $2$--$4\%$ 옮겨 가므로 측정 오차보다 작다. 결정적인 것은 루프를 여는 실험(첫 단 입력에 최종 호르몬을 일정하게 주면 이런 리듬은 사라진다)이나, 주기의 여러 위상에서 호르몬을 짧게 추가했을 때의 반응이다.}
